\documentclass[aspectratio=169]{beamer}
\usetheme{Madrid}
\usecolortheme{default}
\usepackage[utf8]{inputenc}
\usepackage{amsmath}
\usepackage{amsfonts}
\usepackage{amssymb}

\setbeamertemplate{navigation symbols}{}
\setbeamerfont{block body}{size=\small}
\setbeamerfont{alertblock body}{size=\small}

\title{CRISC Domain B -- Impossible Questions}
\subtitle{100 Expert-Level Questions -- Answer on the Next Slide}
\author{CRISC Mastery Series}
\date{\today}

% ------------------------------------------------------------
% \hq{number}{topic}{question}{options}{letter}{explanation}{trap}
% Creates one question frame and one explanation frame.
% (No blank lines are allowed inside the arguments.)
% ------------------------------------------------------------
\newcommand{\hq}[7]{%
  \begin{frame}{Q#1: #2}
    \begin{block}{Question}
      #3
    \end{block}
    \begin{enumerate}[(A)]
      #4
    \end{enumerate}
  \end{frame}
  \begin{frame}{Q#1: Explanation}
    \begin{block}{Correct Answer: #5}
      \begin{itemize}
        #6
      \end{itemize}
    \end{block}
    \begin{alertblock}{Critical Trap}
      #7
    \end{alertblock}
  \end{frame}
}

\begin{document}

\begin{frame}
  \titlepage
\end{frame}

\begin{frame}{Warning}
  \begin{alertblock}{Extreme Difficulty}
    100 questions. Each answer and explanation is on the slide after its question.\\
    What remains needs calculation, probability, dependency analysis and judgment.\\
    Do not be discouraged if you miss most of them.
  \end{alertblock}
\end{frame}

\begin{frame}{Outline}
  \tableofcontents
\end{frame}

% ============================================================
\section{Quantitative Risk Analysis}
% ============================================================

\hq{1}{Control ROI}
{An asset is worth \$2M with an exposure factor of 40\% and an ARO of 0.5. A control reduces the ARO to 0.1 and the exposure factor to 25\% at a cost of \$60K per year. What is the net annual benefit?}
{\item \$350K
 \item \$410K
 \item \$290K
 \item \$50K}
{C}
{\item Before: ALE $=$ \$2M $\times$ 40\% $\times$ 0.5 $=$ \$400K
 \item After: ALE $=$ \$2M $\times$ 25\% $\times$ 0.1 $=$ \$50K
 \item Gross reduction $=$ \$350K
 \item Net benefit $=$ \$350K $-$ \$60K $=$ \$290K}
{\$350K ignores the control cost. \$50K is the residual ALE, not a benefit.}

\hq{2}{Joint Probability Bounds}
{Risk A has a 20\% likelihood and Risk B has a 30\% likelihood. They are known to be positively correlated. Which statement MUST be true about the probability that both occur?}
{\item It equals exactly 6\%
 \item It is at least 50\%
 \item It can exceed 30\%
 \item It is strictly greater than 6\% and at most 20\%}
{D}
{\item Independence gives $0.2 \times 0.3 = 6\%$
 \item Positive correlation pushes the joint probability above 6\%
 \item A joint event can never be more likely than its rarer component
 \item Upper bound $= \min(20\%, 30\%) = 20\%$}
{Correlation changes the joint probability, but it is always capped by the least likely event.}

\hq{3}{Annualizing a Monthly Probability}
{A threat event occurs in any given month with probability 5\%, independently of other months. What is the probability of at least one occurrence in 12 months?}
{\item About 46\%
 \item About 54\%
 \item 5\%
 \item 60\%}
{A}
{\item $P(\text{none in 12 months}) = 0.95^{12} \approx 0.54$
 \item $P(\text{at least one}) = 1 - 0.54 \approx 0.46$
 \item 54\% is the probability of NO event
 \item $12 \times 5\% = 60\%$ ignores overlap of outcomes}
{Probabilities of ``at least one'' do not add linearly. Use the complement.}

\hq{4}{Alert Reliability}
{The probability of a targeted attack this quarter is 10\%. When an attack is under way, the detection alert fires 90\% of the time. When no attack is under way, it fires falsely 20\% of the time. The alert has just fired. What is the probability that an attack is actually under way?}
{\item 90\%
 \item 10\%
 \item About 33\%
 \item 50\%}
{C}
{\item $P(\text{alert}) = 0.9 \times 0.1 + 0.2 \times 0.9 = 0.27$
 \item $P(\text{attack} \mid \text{alert}) = 0.09 / 0.27 = 1/3$
 \item A low base rate makes most alerts false positives
 \item The alert raises the estimate from 10\% to 33\%, not to 90\%}
{Detection accuracy is not the same as the probability that an alert is real. Base rates matter.}

\hq{5}{Three-Point Estimate}
{Analysts estimate a loss at \$100K (optimistic), \$200K (most likely) and \$700K (pessimistic). What is the PERT expected loss?}
{\item \$333K
 \item \$200K
 \item \$267K
 \item \$400K}
{C}
{\item PERT mean $= (O + 4M + P)/6$
 \item $= (100 + 800 + 700)/6 \approx \$267K$
 \item \$333K is the simple average of the three values
 \item \$200K is the mode and ignores the long tail}
{The most likely value is not the expected value when the distribution is skewed.}

\hq{6}{Skewed Confidence Interval}
{A FAIR analysis gives a 90\% confidence interval of \$200K to \$2M with a most likely value of \$500K. Which interpretation is BEST?}
{\item Expected annual loss likely exceeds \$500K because of right skew
 \item Expected annual loss is exactly \$500K
 \item The interval is invalid because it is asymmetric
 \item The loss will certainly fall inside the interval}
{A}
{\item Loss distributions are typically right-skewed (lognormal-like)
 \item The long upper tail pulls the mean above the mode
 \item Budgeting on the mode understates expected loss
 \item A 90\% interval leaves a 10\% chance of falling outside it}
{Most likely $\neq$ expected value. Asymmetric intervals are normal for losses.}

\hq{7}{Aggregating Volatility}
{Two risks each have an annual loss standard deviation of \$300K. Their correlation is 0.5. What is the standard deviation of the combined loss?}
{\item \$424K
 \item \$300K
 \item \$600K
 \item \$520K}
{D}
{\item $\text{Var} = \sigma^2 + \sigma^2 + 2\rho\sigma^2 = 3\sigma^2$
 \item $\text{SD} = 300 \times \sqrt{3} \approx \$520K$
 \item \$424K assumes independence ($\rho = 0$)
 \item \$600K assumes perfect correlation ($\rho = 1$)}
{Standard deviations do not add. Correlation controls how much diversification you get.}

\hq{8}{Common-Cause Failure}
{Two preventive controls each stop an attack 90\% of the time. A design review calls them ``independent layers,'' implying a 1\% chance that both fail. Testing reveals a shared dependency (the same identity provider) whose failure disables both controls 4\% of the time. Otherwise they fail independently. What is the probability that both fail?}
{\item 1\%
 \item 4\%
 \item About 5\%
 \item 10\%}
{C}
{\item $P(\text{both fail}) = 0.04 + 0.96 \times 0.10 \times 0.10$
 \item $= 0.04 + 0.0096 \approx 4.96\%$
 \item Residual risk is about five times the ``independent'' estimate
 \item Defense in depth fails when layers share a failure mode}
{Layered controls are only independent if they share no common dependency.}

\hq{9}{Rule of Three}
{A control is tested 30 times with zero failures. Which statement about its true failure rate is MOST defensible?}
{\item It is 0\%
 \item It is at most 3\%
 \item It is at most 30\%
 \item It is at most roughly 10\% with 95\% confidence}
{D}
{\item With zero failures in $n$ trials, the 95\% upper bound is about $3/n$
 \item $3/30 = 10\%$
 \item Zero observed failures does not prove zero failures
 \item Small samples leave wide uncertainty}
{``No failures observed'' is not ``no failures exist.''}

\hq{10}{FAIR Decomposition}
{Threat event frequency is 20 per year and vulnerability (probability a threat event becomes a loss event) is 15\%. Primary loss per event is \$40K. There is a 30\% chance of a secondary loss of \$200K per event. What is the annualized loss exposure?}
{\item \$120K
 \item \$720K
 \item \$300K
 \item \$2.0M}
{C}
{\item Loss event frequency $= 20 \times 15\% = 3$ per year
 \item Expected loss per event $= 40K + 0.3 \times 200K = \$100K$
 \item ALE $= 3 \times \$100K = \$300K$
 \item \$120K ignores secondary loss; \$720K treats it as certain; \$2.0M ignores vulnerability}
{Use loss event frequency, not threat event frequency, and weight secondary loss by its probability.}

% ============================================================
\section{Risk Evaluation}
% ============================================================

\hq{11}{Ordinal Matrix Flaw}
{An organization multiplies 1--5 likelihood and impact scores in a 5x5 matrix. A risk scored 5x1 and a risk scored 1x5 both receive 5. What is the MOST fundamental flaw in this approach?}
{\item Multiplication always overstates risk
 \item Ordinal ranks are not ratio-scale values, so products can misrank risks with very different expected losses
 \item A 5x5 matrix is too small
 \item Scores above 20 cannot be compared}
{B}
{\item Scale points are ordered labels, not measured quantities
 \item Equal products can hide very different loss profiles
 \item Products of ordinal ranks can reverse true priority
 \item Calibrated quantitative ranges are more defensible}
{A numeric-looking score does not make an ordinal scale quantitative.}

\hq{12}{Aggregate Appetite and Correlation}
{Three risks each have an expected annual loss of \$400K. Aggregate appetite is \$1M and aggregate tolerance is \$1.3M. The risks are believed to be strongly positively correlated. Which position is MOST accurate?}
{\item Above appetite but within tolerance on an expected basis; correlation does not change the expected sum but fattens the tail, so the loss distribution must be examined
 \item Within appetite because every individual risk is below \$1M
 \item Beyond tolerance because correlated risks add up to more than their expected values
 \item Within appetite because diversification lowers the expected aggregate loss}
{A}
{\item Expected losses add: $3 \times \$400K = \$1.2M$
 \item \$1.2M is above appetite (\$1M) and within tolerance (\$1.3M)
 \item Correlation does not change the expected sum
 \item Correlation does increase tail (extreme-loss) exposure}
{Correlation affects the spread of outcomes, not the sum of expected values.}

\hq{13}{Cost Exceeds Benefit}
{A control costing \$400K per year would reduce expected loss by only \$150K per year. Without it, residual risk exceeds appetite but is well within capacity. What should the risk practitioner recommend?}
{\item Implement the control because the risk exceeds appetite
 \item Escalate for formal, documented acceptance by the owner with authority to exceed appetite, while exploring cheaper alternatives
 \item Accept informally to avoid delay
 \item Buy full insurance without further analysis}
{B}
{\item Treatment must be economically justified
 \item Exceeding appetite requires explicit authority to accept
 \item Acceptance must be documented and reviewed
 \item Cheaper responses should still be explored}
{Exceeding appetite does not force a treatment. It forces a decision by the right authority.}

\hq{14}{Mitigation That Backfires}
{After a new control is deployed, the verified residual risk is HIGHER than the original inherent risk. What is the MOST likely explanation?}
{\item An arithmetic error is the only plausible cause
 \item The control introduced secondary risk that outweighs its benefit, so the response must be redesigned
 \item Inherent risk was necessarily overstated
 \item The risk has been transferred to the control owner}
{B}
{\item Controls can create new risks (complexity, new attack surface, new dependencies)
 \item Secondary risk must be assessed alongside residual risk
 \item Net exposure, not control count, is what matters
 \item Redesign or replace the response}
{A control is not automatically net-positive.}

% ============================================================
\section{Business Impact Analysis}
% ============================================================

\hq{15}{Recovery Dependencies}
{System A has an RTO of 4 hours. A's own restoration takes 2 hours and can start only after both System B and System C are available. B takes 3 hours to restore and C takes 2 hours, restored in parallel. Can A meet its RTO?}
{\item Yes, the longest single task is 3 hours
 \item Yes, B and C run in parallel and fit inside 4 hours
 \item No, recovery takes 7 hours
 \item No, recovery takes 5 hours}
{D}
{\item B and C run in parallel, so the dependency stage takes $\max(3, 2) = 3$ hours
 \item A then needs its own 2 hours: $3 + 2 = 5$ hours
 \item 5 hours exceeds the 4-hour RTO
 \item 7 hours would be the fully sequential case}
{Critical path = longest dependency chain, not the longest single task.}

\hq{16}{Probability of Meeting RPO}
{The RPO is 1 hour and backups run every 4 hours. A failure occurs at a uniformly random time. What is the probability that the data loss stays within the RPO?}
{\item 0\%
 \item 25\%
 \item 50\%
 \item 75\%}
{B}
{\item Data loss equals the time since the last backup, uniform from 0 to 4 hours
 \item Loss is within RPO only when it is at most 1 hour
 \item Probability $= 1/4 = 25\%$
 \item Expected loss is 2 hours, twice the RPO}
{A backup interval longer than the RPO means the RPO is usually missed.}

\hq{17}{MTD Consistency}
{A process has an RTO of 4 hours, a work recovery time (WRT) of 6 hours and an MTD of 8 hours. What is the PRIMARY problem?}
{\item Consistent because RTO (4 h) is less than MTD (8 h)
 \item Consistent because WRT is excluded from MTD
 \item MTD should automatically be raised to 20 hours
 \item Inconsistent: RTO plus WRT is 10 hours, which exceeds the 8-hour MTD}
{D}
{\item The business must be fully operating again within the MTD
 \item Requirement: $\text{RTO} + \text{WRT} \leq \text{MTD}$
 \item Here $4 + 6 = 10 > 8$
 \item Shorten RTO or WRT, or revalidate the MTD with the business}
{Meeting the RTO alone does not guarantee meeting the MTD.}

\hq{18}{MTD from Escalating Impact}
{Outage losses are \$10K per hour for the first 4 hours and \$50K per hour afterwards. Management defines MTD as the point where cumulative loss reaches \$400K. What is the MTD?}
{\item 11.2 hours
 \item 8 hours
 \item 4 hours
 \item 40 hours}
{A}
{\item First 4 hours: $4 \times \$10K = \$40K$
 \item Remaining threshold: $\$400K - \$40K = \$360K$
 \item At \$50K per hour: $360 / 50 = 7.2$ hours
 \item MTD $= 4 + 7.2 = 11.2$ hours}
{Impact is not linear. Piece the loss curve together before dividing.}

% ============================================================
\section{Risk Treatment}
% ============================================================

\hq{19}{Treatment Selection}
{A risk has an 80\% likelihood and a \$100K impact. Treatment A costs \$50K and cuts likelihood to 20\%. Treatment B costs \$30K and cuts impact to \$40K. Applying both costs \$80K. Which option has the best net benefit?}
{\item Both treatments
 \item Treatment B only
 \item Neither
 \item Treatment A only}
{B}
{\item Baseline expected loss $= 0.8 \times 100 = \$80K$
 \item A: loss \$20K, net benefit $= 60 - 50 = +\$10K$
 \item B: loss \$32K, net benefit $= 48 - 30 = +\$18K$
 \item Both: loss \$8K, net benefit $= 72 - 80 = -\$8K$}
{Benefits of overlapping treatments are not additive. The second one has diminishing value.}

\hq{20}{Present Value of a Control}
{A control costs \$120K upfront and reduces expected annual loss from \$90K to \$20K for 3 years. The discount rate is 10\%. What is the approximate NPV?}
{\item About $+$\$54K, so approve
 \item About $+$\$90K
 \item About $-$\$50K, so reject
 \item About $+$\$210K}
{A}
{\item Annual benefit $= \$70K$
 \item PV factor for 3 years at 10\% $\approx 2.487$
 \item PV of benefits $\approx \$174K$
 \item NPV $\approx 174 - 120 = +\$54K$}
{\$90K is the undiscounted net. \$210K is gross undiscounted benefit.}

\hq{21}{Insurance Economics}
{There is a 10\% annual chance of a \$2M loss. Insurance covers 80\% of the loss above a \$200K deductible and costs \$60K per year. What is the total expected annual cost with insurance?}
{\item \$56K
 \item \$200K
 \item \$116K
 \item \$260K}
{C}
{\item Retained per event $= 200K + 20\% \times 1.8M = \$560K$
 \item Expected retained loss $= 0.1 \times 560K = \$56K$
 \item Add premium: $56K + 60K = \$116K$
 \item Uninsured expected loss would be \$200K}
{Include the premium and the retained share. Insurance also leaves counterparty and coverage risk.}

% ============================================================
\section{Risk Monitoring}
% ============================================================

\hq{22}{Control Test Results}
{A control runs 250 times per year. A sample of 25 executions shows 5 deviations. The tolerable deviation rate is 5\%. What is the BEST conclusion?}
{\item Effective, because the sample is only 10\% of the population
 \item Effective, because 80\% of items passed
 \item Inconclusive until all 250 executions are tested
 \item Not effective: the 20\% sample deviation rate far exceeds the 5\% tolerable rate, so remediate and expand testing}
{D}
{\item Sample deviation rate $= 5/25 = 20\%$
 \item 20\% is four times the tolerable rate
 \item A passing rate of 80\% is not acceptable here
 \item Full-population testing is not needed to see a clear failure}
{Compare results to the tolerable rate, not to a ``mostly passed'' feeling.}

\hq{23}{Silent Indicators}
{Every KRI has stayed green for three consecutive years, yet the organization suffered two significant incidents. What is the MOST appropriate conclusion?}
{\item Controls are exceptionally effective
 \item Thresholds are likely too loose or indicators insensitive, so back-test and recalibrate
 \item Reduce monitoring frequency to save cost
 \item Retire the KRIs}
{B}
{\item Indicators should have moved before the incidents
 \item Back-test KRIs against past incidents
 \item Tighten thresholds and add leading indicators
 \item Consistent green is a warning sign, not proof of safety}
{A KRI that cannot detect a known failure is not doing its job.}


% ============================================================
\section{Advanced Quantitative Analysis}
% ============================================================

\hq{24}{Poisson Annual Probability}
{Loss events arrive at an average rate of 0.5 per year (Poisson, independent). What is the probability of at least one event in a year?}
{\item 25\%
 \item About 39\%
 \item 50\%
 \item About 63\%}
{B}
{\item $P(\text{none}) = e^{-0.5} \approx 0.607$
 \item $P(\text{at least one}) \approx 0.393$
 \item 63\% would be the answer for a rate of 1.0 per year
 \item 50\% wrongly treats the rate as a probability}
{A rate is not a probability. Use $1 - e^{-\lambda}$.}

\hq{25}{Portfolio Compromise}
{Ten independent systems each have a 10\% annual chance of compromise. What is the probability that at least one is compromised?}
{\item About 35\%
 \item 100\%
 \item About 65\%
 \item 10\%}
{C}
{\item $P(\text{none}) = 0.9^{10} \approx 0.349$
 \item $P(\text{at least one}) \approx 0.651$
 \item 100\% comes from adding $10 \times 10\%$
 \item 35\% is the probability that none are compromised}
{Many small independent risks add up to a large combined risk, but not linearly.}

\hq{26}{Exclusive Scenarios}
{Three mutually exclusive outcomes this year: a 10\% chance of a \$500K loss, a 5\% chance of a \$2M loss, otherwise no loss. What is the expected loss?}
{\item \$150K
 \item \$2.5M
 \item \$250K
 \item \$75K}
{A}
{\item $0.10 \times 500K = \$50K$
 \item $0.05 \times 2M = \$100K$
 \item Sum $= \$150K$
 \item \$2.5M adds the severities without weighting}
{Weight each loss by its own probability.}

\hq{27}{Volatility of One Event}
{A single event has a 20\% chance of causing a \$1M loss and otherwise causes nothing. What is the standard deviation of the loss?}
{\item \$200K
 \item \$800K
 \item \$160K
 \item \$400K}
{D}
{\item Bernoulli loss: $\text{SD} = L\sqrt{p(1-p)}$
 \item $= 1M \times \sqrt{0.2 \times 0.8} = 1M \times 0.4$
 \item \$200K is the expected loss, not the volatility
 \item Rare large losses are very volatile}
{Expected loss is not the same as volatility.}

\hq{28}{Value at Risk}
{Annual loss is normally distributed with mean \$1M and standard deviation \$400K. What is the 95\% one-tailed VaR ($z = 1.645$)?}
{\item About \$1.66M
 \item \$1.00M
 \item About \$1.93M
 \item \$1.40M}
{A}
{\item $\text{VaR} = \mu + z\sigma = 1M + 1.645 \times 0.4M$
 \item $\approx \$1.66M$
 \item \$1.93M is the 99\% figure ($z = 2.326$)
 \item \$1.00M is only the mean}
{VaR is a quantile of the loss distribution, not the average loss.}

\hq{29}{Expected Shortfall}
{Which statement about 95\% expected shortfall (ES) is correct?}
{\item It is the average loss in the worst 5\% of outcomes and is never below the 95\% VaR
 \item It is the average loss across all outcomes
 \item It is the loss level exceeded 5\% of the time
 \item It always equals VaR for normal distributions}
{A}
{\item VaR marks where the tail starts
 \item ES averages the losses beyond that point
 \item So $\text{ES} \geq \text{VaR}$
 \item ES captures tail severity that VaR ignores}
{VaR says where the tail begins. ES says how bad it gets.}

\hq{30}{Monte Carlo Precision}
{A Monte Carlo run of 10,000 independent trials estimates an event probability of 5\%. What is the approximate standard error?}
{\item About 0.22 percentage points
 \item About 0.5 percentage points
 \item About 0.05 percentage points
 \item About 2.2 percentage points}
{A}
{\item $\text{SE} = \sqrt{p(1-p)/n}$
 \item $= \sqrt{0.05 \times 0.95 / 10000} \approx 0.0022$
 \item So the estimate is about $5\% \pm 0.4\%$ at 95\% confidence
 \item More trials shrink error only with the square root of $n$}
{Simulation output is an estimate with its own sampling error.}

\hq{31}{Reverse Conditional}
{2\% of systems are breached each year. Among breached systems, 80\% were unpatched. 20\% of all systems are unpatched. What is the probability that an unpatched system is breached?}
{\item 80\%
 \item 2\%
 \item 8\%
 \item 40\%}
{C}
{\item Bayes: $P(B \mid U) = P(U \mid B)P(B)/P(U)$
 \item $= 0.8 \times 0.02 / 0.20 = 0.08$
 \item 80\% is the reverse conditional probability
 \item Unpatched systems are 4 times riskier than average}
{$P(A \mid B) \neq P(B \mid A)$.}

\hq{32}{Multiplicative Sensitivity}
{$\text{ALE} = \text{TEF} \times \text{vulnerability} \times \text{loss}$. TEF rises 50\%, vulnerability falls 40\%, and loss per event rises 25\%. What happens to ALE?}
{\item It rises 35\%
 \item It rises 10\%
 \item It falls 5\%
 \item It rises 12.5\%}
{D}
{\item Factors multiply: $1.5 \times 0.6 \times 1.25 = 1.125$
 \item Adding the percentages gives $+35\%$, which is wrong
 \item Order of changes does not matter
 \item Run one-at-a-time sensitivities to see drivers}
{Percent changes in multiplicative models compound. They do not add.}

\hq{33}{Perfect Hedge}
{Two risks each have a loss standard deviation of \$300K and a correlation of $-1$. What is the standard deviation of the combined loss?}
{\item \$300K
 \item \$600K
 \item \$0
 \item \$424K}
{C}
{\item $\text{Var} = \sigma^2 + \sigma^2 + 2\rho\sigma^2 = 0$ when $\rho = -1$
 \item Gains in one exactly offset losses in the other
 \item Perfect negative correlation is rare in practice
 \item Never assume it without evidence}
{Diversification benefit depends entirely on correlation.}

\hq{34}{Frequency Versus Severity}
{Scenario X: 4 events per year of \$250K each. Scenario Y: 2 events per year of \$500K each. Both are Poisson with fixed severity and the same \$1M expected annual loss. Which is true?}
{\item Volatility is identical
 \item Y is more volatile (SD about \$707K versus \$500K)
 \item X is more volatile
 \item Y has a lower expected loss}
{B}
{\item Compound Poisson: $\text{Var} = \lambda s^2$
 \item X: $4 \times 250^2 \to \text{SD} = \$500K$
 \item Y: $2 \times 500^2 \to \text{SD} \approx \$707K$
 \item Fewer, larger events means a fatter tail}
{Equal expected loss does not mean equal risk.}

\hq{35}{Premium Loading}
{Expected annual loss is \$200K and an insurer charges \$260K. What is the loading on expected loss?}
{\item 26\%
 \item 23\%
 \item 60\%
 \item 30\%}
{D}
{\item Loading $= (\text{premium} - \text{EL})/\text{EL}$
 \item $= (260 - 200)/200 = 30\%$
 \item 23\% wrongly divides by the premium
 \item Insurance usually costs more than expected loss}
{Measure loading against expected loss, not against the premium.}

\hq{36}{Backing Out Inherent Risk}
{After controls with combined 70\% effectiveness, residual exposure is \$120K. What was the inherent exposure?}
{\item \$360K
 \item \$400K
 \item \$204K
 \item \$171K}
{B}
{\item $\text{Residual} = \text{Inherent} \times (1 - 0.7)$
 \item $\text{Inherent} = 120 / 0.3 = \$400K$
 \item \$171K divides by 0.7 by mistake
 \item \$204K multiplies by 1.7 by mistake}
{Residual is a fraction of inherent, so divide by the remaining fraction.}

\hq{37}{Return Period}
{A loss level has a 1-in-20-year return period. What is the probability it is exceeded at least once in 10 years?}
{\item About 40\%
 \item 50\%
 \item 5\%
 \item 10\%}
{A}
{\item Annual probability $= 5\%$
 \item $P(\text{none in 10 years}) = 0.95^{10} \approx 0.60$
 \item $P(\text{at least one}) \approx 40\%$
 \item 50\% would require 14 years}
{A return period is not a schedule. Events can cluster or never occur.}

\hq{38}{Century Event}
{What is the probability of at least one 1-in-100-year event in a 30-year period?}
{\item About 26\%
 \item 1\%
 \item 30\%
 \item 74\%}
{A}
{\item $P(\text{none}) = 0.99^{30} \approx 0.74$
 \item $P(\text{at least one}) \approx 0.26$
 \item 74\% is the probability of none
 \item 30\% adds probabilities linearly}
{Use the complement for at-least-one probabilities.}

\hq{39}{Perpetual Benefit}
{A control costs \$600K upfront and permanently cuts expected annual loss by \$60K. The discount rate is 8\%. What is the NPV?}
{\item $+$\$750K
 \item $-$\$540K
 \item $-$\$150K
 \item $+$\$150K}
{D}
{\item PV of a perpetuity $= 60 / 0.08 = \$750K$
 \item $\text{NPV} = 750 - 600 = +\$150K$
 \item \$750K is the benefit, not the NPV
 \item Payback is 10 years, but NPV is still positive}
{Value permanent benefits as a perpetuity, then subtract cost.}

\hq{40}{Calibration}
{Experts give 90\% confidence ranges, but actual values fall inside the ranges only 50\% of the time. What is the diagnosis?}
{\item Overconfidence: ranges are too narrow, so apply calibration training
 \item Normal behavior for 90\% ranges
 \item Underconfidence: ranges are too wide
 \item Only a data quality problem}
{A}
{\item A well-calibrated 90\% range captures truth about 90\% of the time
 \item 50\% capture means ranges are far too tight
 \item Calibration training and feedback improve this
 \item Widen ranges or use structured elicitation}
{Most experts are overconfident. Test calibration.}

\hq{41}{Budgeted Portfolio}
{Budget is \$250K. Independent, additive options: A costs \$50K and avoids \$120K; B costs \$200K and avoids \$400K; C costs \$20K and avoids \$35K. Which portfolio maximizes net benefit?}
{\item A and B (net \$270K)
 \item A and C (net \$85K)
 \item A, B and C (net \$285K)
 \item B and C (net \$215K)}
{A}
{\item Net: A $= +70$, B $= +200$, C $= +15$
 \item A + B costs \$250K, net \$270K
 \item A + B + C costs \$270K and breaks the budget
 \item B + C nets only \$215K}
{Highest ROI alone is not the same as maximum value under a budget.}

\hq{42}{Two-Stage Detection}
{Stage 1 detection catches 70\% of intrusions. Whatever passes is examined by stage 2, which catches 50\% of what it sees (independent of stage 1). What fraction of intrusions is missed by both?}
{\item 15\%
 \item 0\%
 \item 35\%
 \item 20\%}
{A}
{\item Missed by stage 1 $= 30\%$
 \item Missed by stage 2 as well $= 0.3 \times 0.5 = 15\%$
 \item 20\% would wrongly average the catch rates
 \item Detection layers multiply the misses}
{Combined misses multiply. Catch rates do not add.}

\hq{43}{Weighted Scenario Mean}
{Scenario probabilities and losses: 60\% at \$50K, 30\% at \$200K, 10\% at \$1M. What is the expected loss?}
{\item \$200K
 \item About \$417K
 \item \$190K
 \item \$1.25M}
{C}
{\item $0.6 \times 50 = 30$
 \item $0.3 \times 200 = 60$
 \item $0.1 \times 1000 = 100$; total $\$190K$
 \item \$417K is the unweighted average of the three losses}
{Always weight outcomes by probability.}

\hq{44}{Multi-Year Poisson}
{Loss events arrive at 0.5 per year (Poisson). What is the probability of at least one event in 3 years?}
{\item About 22\%
 \item 50\%
 \item About 78\%
 \item About 85\%}
{C}
{\item Expected events in 3 years $= 1.5$
 \item $P(\text{none}) = e^{-1.5} \approx 0.223$
 \item $P(\text{at least one}) \approx 0.777$
 \item 22\% is the probability of none}
{Scale the rate by time, then use the Poisson complement.}

% ============================================================
\section{Advanced Risk Evaluation}
% ============================================================

\hq{45}{Usable Appetite Statement}
{Which risk appetite statement is MOST usable for decisions?}
{\item The organization will follow industry standards
 \item Expected annual cyber loss will not exceed \$2M, and no scenario will exceed \$5M at 95\% confidence
 \item The organization has a low appetite for cyber risk
 \item The organization will manage cyber risk responsibly}
{B}
{\item It is measurable
 \item It sets both an aggregate and a tail limit
 \item It names a confidence level
 \item Vague words cannot trigger escalation}
{If it cannot be measured, it cannot be breached.}

\hq{46}{Beyond Tolerance}
{Capacity is \$10M, tolerance is \$6M and appetite is \$4M. Aggregate exposure is now \$7M. What is the BEST description?}
{\item Within appetite because it is below capacity
 \item Beyond tolerance but within capacity: escalate and either reduce exposure or formally approve it
 \item Capacity exceeded: stop operations
 \item Acceptable because it is below capacity}
{B}
{\item Ordering is appetite $<$ tolerance $<$ capacity
 \item \$7M exceeds tolerance (\$6M)
 \item It is still under capacity (\$10M)
 \item Tolerance breach triggers escalation and a plan}
{Capacity is the survival limit, not the acceptance limit.}

\hq{47}{Central Tendency}
{70\% of assessed risks cluster at ``medium on a 5-point scale. What is the most likely problem?}
{\item The analysts are too senior
 \item Central tendency bias: use anchored criteria and calibrate scorers
 \item The matrix is too large
 \item The portfolio is genuinely balanced}
{B}
{\item Raters avoid extreme scores
 \item Vague scale definitions encourage hedging
 \item Anchored examples reduce it
 \item Compare scorer distributions to detect it}
{A suspiciously uniform distribution signals a scoring problem.}

\hq{48}{Same Mean, Different Tail}
{Risk X has expected loss \$100K and SD \$50K. Risk Y has expected loss \$100K and SD \$400K. The board has low tolerance for loss volatility. Which should be prioritized?}
{\item Y: equal expected loss hides a far larger tail
 \item Equal priority
 \item X: it is more certain
 \item Neither: expected losses are equal}
{A}
{\item Y can produce losses many times its mean
 \item Low tolerance for volatility means tail matters
 \item Ranking by mean alone ignores this
 \item Evaluate distributions, not only point values}
{Expected value ranks risks only if volatility is irrelevant.}

\hq{49}{Concentration}
{Twelve risks rated ``medium all depend on the same vendor. What is the BEST interpretation?}
{\item Low risk, since medium is acceptable
 \item A concentration risk to evaluate in aggregate, which may exceed appetite even though each item is medium
 \item A procurement issue only
 \item Twelve independent mediums, so no action}
{B}
{\item A shared dependency creates correlated failure
 \item The portfolio exposure is greater than any single rating
 \item Aggregate by dependency, not only by item
 \item Registers often hide concentration}
{Twelve mediums with one root cause is one large risk.}

\hq{50}{Silo Aggregation}
{Five departments each accept a \$200K risk (each limit is \$250K). All depend on one shared service. Enterprise aggregate appetite is \$600K. What is the issue?}
{\item Appetite applies per department only
 \item The exposure is only \$200K because it is the same dependency
 \item None: each is within its limit
 \item Correlated exposure could total \$1M, above enterprise appetite, so acceptances must be aggregated}
{D}
{\item $5 \times \$200K = \$1M$ if the shared service fails
 \item Local limits do not guarantee enterprise limits
 \item Acceptance needs an enterprise view
 \item Escalate or treat the shared dependency}
{Local approval does not equal enterprise acceptance.}

\hq{51}{Residual-Only Register}
{A risk register records only residual ratings. What is the main weakness?}
{\item Residual ratings are always inaccurate
 \item Registers must list threats only
 \item Inherent risk never matters
 \item Hidden dependence on controls: if a control fails, the true exposure is invisible}
{D}
{\item Inherent rating shows exposure without controls
 \item It shows how much each control is carrying
 \item Control failure analysis needs both
 \item Treatment prioritization needs both}
{Residual alone hides control dependence.}

\hq{52}{Velocity}
{Two risks both score 12. Risk A can reach full impact within minutes with no warning. Risk B develops over months with clear leading indicators. Where is stronger preventive investment most needed?}
{\item Risk B, because indicators are costly
 \item Neither: scores are equal
 \item Risk B, because it lasts longer
 \item Risk A, because short velocity leaves no time to respond after detection}
{D}
{\item Velocity is time from trigger to impact
 \item A has no response window
 \item B can be managed with monitoring and response
 \item Equal scores can need very different strategies}
{Same score, different urgency.}

\hq{53}{Rank Reversal}
{Impact levels are logarithmic: level 2 $=$ \$100K and level 5 $=$ \$100M. Risk P has likelihood level 5 (50\%) and impact 2. Risk Q has likelihood level 2 (5\%) and impact 5. Both score 10. What is the real difference?}
{\item Q has about 100 times the expected loss (\$5M versus \$50K)
 \item Q has 10 times the expected loss
 \item P has the larger expected loss
 \item They are equal}
{A}
{\item P: $0.5 \times 100K = \$50K$
 \item Q: $0.05 \times 100M = \$5M$
 \item Equal scores hide a 100x gap
 \item Multiplying ordinal levels ignores the underlying scales}
{Matrix scores are not expected losses.}

\hq{54}{Appetite Test}
{Appetite: no scenario will exceed \$5M at 95\% confidence. Simulation shows a 95th percentile of \$4.8M and a 99th percentile of \$9M. What is the BEST conclusion?}
{\item Compliant, so no further discussion
 \item Compliant with the statement as written, but the 99\% tail deserves board attention and a review of the statement
 \item Non-compliant because the 99th percentile exceeds \$5M
 \item Cannot be tested}
{B}
{\item 95th percentile \$4.8M is under \$5M
 \item The statement fixed the confidence level at 95\%
 \item The 99\% tail is much larger
 \item Choose confidence levels deliberately}
{Passing a test is not the same as being safe.}

% ============================================================
\section{Advanced Business Impact and Recovery}
% ============================================================

\hq{55}{Availability Budget}
{A service has a 99.9\% availability target measured over a 30-day month. How much downtime does it permit?}
{\item About 7.2 hours
 \item About 4.3 minutes
 \item About 43 minutes
 \item About 8.8 hours}
{C}
{\item 30 days $= 43{,}200$ minutes
 \item $0.1\% \times 43{,}200 = 43.2$ minutes
 \item 7.2 hours would be 99\%
 \item 8.8 hours per year is the figure for 99.9\% annually}
{Check whether the target period is monthly or yearly.}

\hq{56}{Series Availability}
{Three components in series each have 99.9\% availability and all must work. What is the combined availability?}
{\item 99.9\%
 \item 99.97\%
 \item About 99.7\%
 \item 99.0\%}
{C}
{\item $0.999^3 \approx 0.997$
 \item Series dependencies multiply
 \item Each added component lowers availability
 \item The weakest link is not the whole story}
{In series, availability multiplies down.}

\hq{57}{Parallel Availability}
{Two independent redundant components each have 99\% availability and either is sufficient. What is the combined availability?}
{\item 99.99\%
 \item 99.9\%
 \item 99\%
 \item 98\%}
{A}
{\item $1 - 0.01^2 = 0.9999$
 \item Both must fail for an outage
 \item Assumes independence and no common cause
 \item Shared power or software erodes the benefit}
{Redundancy helps only if failures are independent.}

\hq{58}{MTBF and MTTR}
{A system has an MTBF of 1,000 hours and an MTTR of 4 hours. What is its steady-state availability?}
{\item 96.00\%
 \item 99.96\%
 \item About 99.60\%
 \item 99.00\%}
{C}
{\item $A = \text{MTBF} / (\text{MTBF} + \text{MTTR})$
 \item $= 1000 / 1004 \approx 0.9960$
 \item Repair time matters as much as failure rate
 \item 99.96\% would need an MTTR of 0.4 hours}
{Availability depends on both how often and how long.}

\hq{59}{Annual Downtime}
{Using MTBF 1,000 hours and MTTR 4 hours, what is the expected annual downtime?}
{\item 4 hours
 \item About 87.6 hours
 \item About 35 hours
 \item About 8.8 hours}
{C}
{\item Unavailability $= 4 / 1004 \approx 0.398\%$
 \item $8760 \times 0.00398 \approx 34.9$ hours
 \item Each failure costs 4 hours; there are about 8.7 failures per year
 \item A small unavailability still adds up}
{Convert percentages to hours per year before judging them.}

\hq{60}{Replication Lag}
{The RPO is 5 minutes. Asynchronous replication typically lags 15 minutes. What is the conclusion?}
{\item RPO is not met: lag exceeds RPO, so reduce lag or use synchronous replication
 \item Nothing can be concluded
 \item RPO is met because RTO is unaffected
 \item RPO is met because replication exists}
{A}
{\item Data in the lag window is lost at failover
 \item 15 minutes is greater than 5 minutes
 \item Replication is not the same as meeting RPO
 \item Choose technology to fit the objective}
{Having replication does not mean meeting the RPO.}

\hq{61}{Sequential Recovery}
{MTD is 24 hours. Recovery is strictly sequential: network 6 h, database 8 h, application 6 h, validation 6 h. Can the MTD be met?}
{\item No: the total is 32 hours
 \item No: total is 26 hours, which exceeds the MTD by 2 hours
 \item Yes: the total is 20 hours
 \item Yes: the longest step is 8 hours}
{B}
{\item $6 + 8 + 6 + 6 = 26$ hours
 \item Sequential steps add. They do not overlap
 \item Parallelize or pre-stage to save time
 \item Revisit MTD only if the business agrees}
{Sum sequential steps. Do not look at the longest one.}

\hq{62}{Hot Site Economics}
{Downtime costs \$80K per hour. A hot site costs \$500K per year and cuts recovery from 24 hours to 2 hours. A disaster has a 5\% annual probability. What is the conclusion?}
{\item Justified: avoided loss per disaster is \$1.76M, which exceeds \$500K
 \item Not justified on cost alone: expected benefit is about \$88K versus \$500K
 \item Break-even at \$500K
 \item Justified because 24 hours is too long}
{B}
{\item Hours saved $= 22$
 \item Loss avoided per disaster $= 22 \times 80K = \$1.76M$
 \item Expected annual benefit $= 0.05 \times 1.76M = \$88K$
 \item Other drivers (regulation, reputation) may still justify it}
{Weight avoided loss by probability before comparing with annual cost.}

\hq{63}{Recovery Capacity}
{Ten systems each have a 4-hour RTO. The team can restore 2 systems at a time and each wave takes 3 hours. How many systems can meet the RTO?}
{\item 6
 \item 4
 \item 10
 \item 2}
{D}
{\item Wave 1 finishes at 3 hours (2 systems)
 \item Wave 2 finishes at 6 hours, which is over the RTO
 \item Equal RTOs do not guarantee capacity
 \item Capacity planning must match the RTOs}
{RTOs are only achievable if recovery capacity exists.}

\hq{64}{Test Scaling}
{A DR test restored a system in 6 hours. Data has grown 50\% since, and restore time scales linearly with data. What restore time should planners expect?}
{\item About 9 hours
 \item About 7.5 hours
 \item 6 hours
 \item About 12 hours}
{A}
{\item $6 \times 1.5 = 9$ hours
 \item Old test results expire as data grows
 \item RTO of 8 hours would now fail
 \item Retest after material change}
{A passed DR test has a shelf life.}

\hq{65}{Nightly Plus Logs}
{Full backups run nightly and transaction logs are shipped every 15 minutes. What RPO is realistically achievable?}
{\item 12 hours
 \item Zero
 \item About 15 minutes
 \item 24 hours}
{C}
{\item Restore the full backup, then replay logs
 \item Loss is bounded by the log interval
 \item Zero would need synchronous replication
 \item Log shipping itself must be monitored}
{RPO is set by the most frequent recovery point, not the slowest.}

% ============================================================
\section{Advanced Risk Treatment}
% ============================================================

\hq{66}{Break-Even Likelihood}
{A \$90K per year control fully removes a \$600K loss. What annual likelihood makes the control break even?}
{\item 54\%
 \item 90\%
 \item 15\%
 \item 6.7\%}
{C}
{\item Break-even when $p \times 600K = 90K$
 \item $p = 90/600 = 15\%$
 \item If the likelihood is lower, the control loses money
 \item Compare with realistic likelihood ranges}
{Break-even analysis shows how wrong an estimate can be.}

\hq{67}{ARO Reduction}
{A control costs \$90K per year and reduces ARO from 0.3 to 0.1. SLE is \$500K. What is the net annual benefit?}
{\item $-$\$10K
 \item $+$\$60K
 \item $+$\$10K
 \item $+$\$100K}
{C}
{\item Reduction $= (0.3 - 0.1) \times 500K = \$100K$
 \item Net $= 100 - 90 = +\$10K$
 \item \$100K is the gross benefit
 \item Thin margins make estimates critical}
{Net benefit subtracts the cost from the ALE reduction.}

\hq{68}{ROSI}
{A control reduces ALE by \$400K and costs \$160K. What is the return on security investment?}
{\item 250\%
 \item 150\%
 \item 40\%
 \item 60\%}
{B}
{\item $\text{ROSI} = (\text{reduction} - \text{cost}) / \text{cost}$
 \item $= (400 - 160)/160 = 1.5$
 \item 250\% divides gross benefit by cost
 \item 40\% is cost divided by benefit}
{ROSI uses net benefit, not gross benefit.}

\hq{69}{Stacked Controls}
{Three independent layers each remove half of the remaining risk. What fraction of inherent risk remains?}
{\item 12.5\%
 \item 25\%
 \item 0\%
 \item 50\%}
{A}
{\item $0.5^3 = 0.125$
 \item Each layer works on what the last one left
 \item Risk never reaches zero
 \item Assumes independence, which is rarely perfect}
{Layers compound multiplicatively.}

\hq{70}{Prevent Plus Detect}
{Inherent: 10 events per year at \$100K each. A preventive control stops 80\% of events. A detective control halves the impact of remaining events. What is the residual annual loss?}
{\item \$300K
 \item \$100K
 \item \$500K
 \item \$200K}
{B}
{\item Inherent $= 10 \times 100K = \$1M$
 \item After prevention: $2 \times 100K = \$200K$
 \item After detection: $2 \times 50K = \$100K$
 \item \$200K counts only prevention}
{Preventive controls cut frequency, detective ones cut impact. Apply each to the right term.}

\hq{71}{Avoidance Economics}
{A company exits a product line earning \$2M profit per year to avoid \$1.5M of expected annual loss. What is the likely effect?}
{\item Net gain of \$3.5M
 \item Net gain of \$1.5M
 \item No effect
 \item Net loss of about \$0.5M per year unless other costs change}
{D}
{\item Value lost $= \$2M$
 \item Loss avoided $= \$1.5M$
 \item Net $= -\$0.5M$ per year
 \item Treat avoidance as a business decision, not a free action}
{Avoiding a risk also avoids its rewards.}

\hq{72}{SLA Credits}
{A cloud SLA gives service credits capped at one month of fees (\$20K). A multi-day outage would cost \$500K. Does the SLA transfer the risk?}
{\item Yes: the provider bears the loss
 \item No: credits are a small offset and the organization retains most of the loss
 \item Yes: credits are insurance
 \item Only for outages under one day}
{B}
{\item Credits cap at \$20K
 \item \$480K or more stays with the organization
 \item Contractual remedies differ from risk transfer
 \item Consider insurance or resilience}
{An SLA is a performance promise, not an indemnity.}

\hq{73}{Interim Exposure}
{A risk has a 24\% annual probability of a \$1M loss. A treatment takes 6 months to deploy. What is the approximate expected loss during rollout?}
{\item \$240K
 \item \$1M
 \item \$60K
 \item About \$120K}
{D}
{\item Annual expected loss $= 0.24 \times 1M = \$240K$
 \item Six months is half a year: $\$120K$
 \item The risk stays live during implementation
 \item Consider an interim control}
{Delay has a quantifiable price.}

\hq{74}{Insurance Exclusions}
{A cyber policy excludes losses from known vulnerabilities left unpatched for over 30 days. The organization has many such systems. What is the implication?}
{\item The premium makes the risk acceptable
 \item Exclusions do not matter after a claim
 \item The risk is fully transferred
 \item Transfer is partly void: the risk is effectively retained and must be treated by patching}
{D}
{\item Coverage depends on policy conditions
 \item Unpatched systems fall outside cover
 \item Retained risk equals the excluded loss
 \item Review exclusions before relying on transfer}
{Insurance transfers only what the policy actually covers.}

\hq{75}{Cost per Dollar Reduced}
{Control A costs \$100K and reduces \$150K of risk. B costs \$40K and reduces \$80K. C costs \$200K and reduces \$260K. Which is most cost-effective?}
{\item C (\$0.77)
 \item B (\$0.50 per \$1 reduced)
 \item A (\$0.67)
 \item All are equal}
{B}
{\item A: $100/150 = 0.67$
 \item B: $40/80 = 0.50$
 \item C: $200/260 = 0.77$
 \item Lowest cost per dollar reduced wins on efficiency}
{Absolute savings and efficiency can rank differently.}

\hq{76}{Shared Hazard}
{Two data centers 10 km apart sit on the same flood plain and the same power grid. The design calls them geographically redundant. What is the concern?}
{\item Common-cause exposure weakens the redundancy; assess shared hazards
 \item Floods are not a risk for data centers
 \item None: two sites always means independence
 \item Distance is the only factor}
{A}
{\item Redundancy requires independent failure modes
 \item One flood or grid event could hit both
 \item Check power, network and hazard overlap
 \item Independence must be demonstrated}
{Redundancy without independence is an illusion.}

\hq{77}{Overdue Treatment}
{A treatment plan is 6 months overdue and the related KRI is worsening. The risk was accepted on the assumption of timely treatment. What is the BEST action?}
{\item Extend the deadline silently
 \item Wait for completion
 \item Reassess residual risk now and escalate: the approved risk position no longer holds
 \item Close the risk}
{C}
{\item Acceptance depended on the plan
 \item Delay increases exposure
 \item The approver must know
 \item A new decision is required}
{A risk decision based on a plan expires when the plan slips.}

% ============================================================
\section{Advanced Risk Monitoring}
% ============================================================

\hq{78}{Control Indicator}
{The percentage of privileged-account reviews completed on time is BEST classified as a:}
{\item Key goal indicator
 \item Key risk indicator (KRI)
 \item Key control indicator (KCI)
 \item Key performance indicator for revenue}
{C}
{\item It shows whether a control is operating
 \item KRIs show changes in exposure
 \item Controls and risk are monitored separately
 \item A falling KCI may later move a KRI}
{KCI measures the control. KRI measures the risk.}

\hq{79}{Zero Deviations}
{A control was tested with zero deviations in 59 sampled items. What can be supported statistically?}
{\item 95\% confidence that the rate is at most 1\%
 \item Nothing can be supported
 \item About 95\% confidence that the deviation rate is at most roughly 5\%
 \item The deviation rate is 0\%}
{C}
{\item $P(\text{0 of 59 if rate} = 5\%) = 0.95^{59} \approx 0.05$
 \item So rate above 5\% is unlikely at 95\% confidence
 \item Zero observed does not mean zero true rate
 \item Tighter bounds need larger samples}
{Sample results bound the rate. They do not prove it.}

\hq{80}{False Alarm Rate}
{A weekly KRI is approximately normal and the alert threshold is the mean plus 2 standard deviations (one-tailed probability 2.3\%). About how many false alerts per year (52 weeks) should be expected when nothing has changed?}
{\item About 5.2
 \item About 2.6
 \item About 0.1
 \item About 1.2}
{D}
{\item $0.0228 \times 52 \approx 1.2$
 \item 2.6 would use a 5\% tail
 \item A 3 SD threshold gives about 0.07 per year
 \item Thresholds must balance sensitivity and noise}
{Every threshold has a false-alarm rate.}

\hq{81}{Compounding Trend}
{A KRI currently at 50 is growing 8\% per week. The threshold is 100. About how long until it breaches?}
{\item About 9 weeks
 \item About 25 weeks
 \item About 12.5 weeks
 \item About 6 weeks}
{A}
{\item $50 \times 1.08^n = 100$
 \item $n = \ln 2 / \ln 1.08 \approx 9.0$
 \item 12.5 weeks wrongly assumes linear growth (+4 per week)
 \item Compounding shortens time to breach}
{Percentage growth compounds. Do not extrapolate linearly.}

\hq{82}{Decaying Effectiveness}
{A control starts 90\% effective. Without maintenance its effectiveness decays by 10\% (relative) each quarter. What is it after 4 quarters?}
{\item About 59\%
 \item About 54\%
 \item About 81\%
 \item 50\%}
{A}
{\item $0.9 \times 0.9^4 = 0.9 \times 0.656$
 \item $\approx 59\%$
 \item 54\% subtracts 36 points linearly
 \item Unmaintained controls drift quietly}
{Controls degrade unless tested and maintained.}

\hq{83}{Self-Reported Metrics}
{A KRI is compiled by the same team it measures, with no validation. What is the main weakness?}
{\item The KRI is too technical
 \item No weakness: teams know best
 \item Reporting bias: use independent or system-sourced data and validate
 \item The KRI is too frequent}
{C}
{\item Teams have incentives to look good
 \item Unverified data cannot support escalation
 \item Automation reduces manual bias
 \item Spot checks add assurance}
{Whoever is measured should not control the measurement.}

\hq{84}{Self-Review Threat}
{A risk practitioner designed a control and is now asked to test its operating effectiveness. What is the concern?}
{\item Self-review threat to objectivity: use an independent tester
 \item The test will take too long
 \item None: the designer knows the control best
 \item The control will be too strict}
{A}
{\item Designers tend to confirm their own work
 \item Independence strengthens assurance
 \item Walkthroughs can still involve the designer
 \item Testing should be done by someone else}
{Do not mark your own homework.}

\hq{85}{Measurement Change}
{A KRI rose for three weeks after a new tool changed how events are counted. What is the BEST action?}
{\item Escalate immediately as a risk increase
 \item Double the threshold
 \item Validate and restate the baseline before concluding risk increased
 \item Ignore the KRI}
{C}
{\item The change may be an artifact of measurement
 \item Compare old and new methods on overlapping data
 \item Document the change and recalibrate
 \item Escalate only if the underlying risk moved}
{A moved number is not always a moved risk.}

\hq{86}{CSA Bias}
{Control self-assessments rate 95\% of controls effective, but independent tests find only 70\% effective. What is the BEST response?}
{\item Trust the self-assessment
 \item Cancel independent tests
 \item Average the two and stop
 \item Treat self-assessment as biased: add independent sampling and recalibrate}
{D}
{\item Gap of 25 points indicates optimism
 \item Self-assessment is useful for engagement but not assurance
 \item Independent evidence should override
 \item Investigate which controls diverge}
{Self-assessment needs independent verification.}

\hq{87}{Leading Indicator}
{Which is the BEST leading KRI for ransomware exposure?}
{\item Percentage of critical systems missing patches older than 30 days
 \item Number of ransomware incidents last year
 \item Total security budget
 \item Number of audit findings closed}
{A}
{\item It moves before an incident occurs
 \item Incident counts are lagging
 \item Budget does not measure exposure
 \item Closed findings measure activity, not exposure}
{A leading indicator changes before the loss does.}

\hq{88}{Exception Aging}
{40\% of risk exceptions are past their expiry date and have been auto-renewed. What does this indicate?}
{\item A permanent acceptance policy
 \item An acceptance governance failure: require re-approval and escalate
 \item Nothing, since exceptions are low risk
 \item Efficient administration}
{B}
{\item Exceptions must be time-bound
 \item Auto-renewal bypasses re-evaluation
 \item Conditions may have changed
 \item Old exceptions accumulate hidden risk}
{Exceptions are not permanent.}

\hq{89}{Sample Size}
{Control A is tested with 25 items from a population of 25,000. Control B is tested with 25 items from a population of 250. Both show zero deviations. Which gives stronger statistical assurance?}
{\item A, because its population is larger
 \item Neither provides any assurance
 \item About the same: absolute sample size drives confidence, not the percentage of the population
 \item B, because it covers 10\% of the population}
{C}
{\item $0.95^{25} \approx 0.28$ for both on the same logic
 \item Confidence depends mainly on $n$ and the deviation rate
 \item Coverage percentage is a poor proxy
 \item Larger samples tighten the bound}
{Percentage of population sampled is not the measure of assurance.}

% ============================================================
\section{Advanced Identification and Governance}
% ============================================================

\hq{90}{Risk Statement Quality}
{Which is the BEST-formed risk statement?}
{\item Because privileged access is not reviewed, a departing administrator could exfiltrate customer data, causing regulatory fines and customer loss
 \item Hackers
 \item Weak access control
 \item Data breach}
{A}
{\item It states cause, event and consequence
 \item ``Weak access control is only a cause
 \item ``Data breach is only an event
 \item ``Hackers is only a threat source}
{A risk statement must link cause, event and business consequence.}

\hq{91}{Pre-Mortem}
{Which technique BEST counters overconfidence in a project plan?}
{\item An annual survey
 \item Pre-mortem: assume the project has failed and explain why
 \item Benchmarking only
 \item Consensus approval}
{B}
{\item It gives permission to voice doubts
 \item It surfaces risks that optimism hides
 \item Consensus can reinforce groupthink
 \item Run it before commitment}
{Imagine failure to find the risks success hides.}

\hq{92}{Fourth Parties}
{A critical vendor relies on an undisclosed subcontractor. Assessing only the vendor misses what?}
{\item Only legal risk
 \item Only pricing risk
 \item Nothing: only the contract signatory matters
 \item Concentration and flow-down risk: assess critical subcontractors and require contractual flow-down}
{D}
{\item Subcontractors may share one provider
 \item Controls may not flow down
 \item Visibility needs contract rights
 \item Critical chains need mapping}
{Your vendor\textquotesingle s vendor is still your risk.}

\hq{93}{Diluted Ownership}
{A significant risk is owned by a committee of five managers. What is the problem?}
{\item The risk is too large
 \item Diluted accountability: assign a single accountable owner
 \item The committee is too small
 \item None: shared ownership is stronger}
{B}
{\item Everyone owns it, so no one acts
 \item One accountable owner can be held to deadlines
 \item Others can be responsible or consulted
 \item Ownership drives follow-through}
{Shared ownership often means no ownership.}

\hq{94}{Regulatory Lead Time}
{A regulation passes with an 18-month effective date. When should the risk assessment of the compliance gap occur?}
{\item At 17 months
 \item Only after a penalty
 \item After the effective date
 \item Now: lead time is the window to design and fund treatment}
{D}
{\item Late assessment compresses options
 \item Treatment may need budget cycles
 \item Early view lets you plan staging
 \item Reassess as guidance evolves}
{Lead time is only useful if you start at the beginning.}

\hq{95}{Incomplete Scenario}
{A scenario lists a threat and a vulnerability but no affected business objective. What is wrong?}
{\item It is incomplete: without consequence for objectives it cannot be prioritized
 \item It is already quantified
 \item Nothing: threat and vulnerability suffice
 \item It is too detailed}
{A}
{\item Risk is about effect on objectives
 \item Impact is required for ranking
 \item Link the scenario to a process and loss type
 \item Without it, treatment cannot be justified}
{No objective at stake means no risk to rank.}

\hq{96}{Healthy Risk Culture}
{Which metric BEST indicates a healthy risk culture?}
{\item Number of policies published
 \item Number of audits scheduled
 \item Number of training slides
 \item Trended time from detecting an incident or near-miss to escalation}
{D}
{\item It measures behavior, not paperwork
 \item Short, steady times suggest people raise bad news early
 \item Rising times signal hesitation
 \item Pair with near-miss reporting rates}
{Culture shows up in how fast bad news travels.}

\hq{97}{Limit Cascade}
{Business unit risk limits sum to \$5M, but the enterprise appetite is \$3M. When is this acceptable?}
{\item Never: limits must sum to the enterprise figure
 \item Always: units are independent
 \item Only if the CFO agrees verbally
 \item Only if units are unlikely to reach their limits at the same time, which needs correlation analysis}
{D}
{\item Limits can exceed appetite if peaks do not coincide
 \item Correlated stress can fill all limits together
 \item Test the cascade with stress scenarios
 \item Document the diversification assumption}
{Sum of limits versus appetite depends on correlation.}

\hq{98}{Cascading Scenario}
{A cloud outage stops payments, which triggers customer attrition, which triggers covenant breaches. How should this be analyzed?}
{\item As one cascading scenario, evaluating the whole chain
 \item As four separate small risks
 \item Only the last consequence
 \item Only the first event}
{A}
{\item Dependencies create sequences
 \item Separate scoring understates the chain
 \item Model triggers and timing
 \item Identify where the chain can be broken}
{Chains of events must be assessed as chains.}

\hq{99}{Fourth Parties}
{A critical vendor relies on an undisclosed subcontractor. Assessing only the vendor misses what?}
{\item Only pricing risk
 \item Concentration and flow-down risk: assess critical subcontractors and require contractual flow-down
 \item Only legal risk
 \item Nothing: only the contract signatory matters}
{B}
{\item Subcontractors may share one provider
 \item Controls may not flow down
 \item Visibility needs contract rights
 \item Critical chains need mapping}
{The vendor\textquotesingle s vendor is still your risk.}

\hq{100}{Anchoring}
{Estimators are told ``last year\textquotesingle s loss was \$1M before giving their own estimates. Which bias is most likely, and what is the fix?}
{\item Optimism bias: reward higher estimates
 \item Availability bias: remove history entirely
 \item No bias: history is objective
 \item Anchoring: elicit independent estimates first, then discuss}
{D}
{\item The number pulls estimates toward it
 \item Independent first-round estimates reduce this
 \item Compare spreads before consensus
 \item Use history as evidence, not as a prompt}
{The first number heard shapes every number after it.}

\begin{frame}{Key Takeaways for Domain B}
  \begin{block}{Summary}
    \begin{itemize}
      \item \textbf{Analysis:} Use complements, base rates, correlation, skew and tails. Percentages compound.
      \item \textbf{Evaluation:} Aggregate across silos. Ordinal scores and means can mislead.
      \item \textbf{BIA:} Critical path, capacity, availability math, and RTO $+$ WRT $\leq$ MTD.
      \item \textbf{Treatment:} Net, discounted benefit. Layers need independence. Transfer needs coverage.
      \item \textbf{Monitoring:} Calibrate thresholds, verify independently, watch for decay.
    \end{itemize}
  \end{block}
\end{frame}

\begin{frame}{Thank You}
  \begin{center}
    {\Huge \textbf{Good Luck on Your CRISC Exam!}}

    \vspace{1cm}

    Remember: the ISACA mindset is about\\
    understanding why a question is asked,\\
    not just what is technically correct.
  \end{center}
\end{frame}

\end{document}
